\documentclass[11pt,a4paper]{article}

% ---------- Packages ----------
\usepackage[margin=1in]{geometry}
\usepackage{amsmath, amssymb, amsthm, mathtools}
\usepackage{microtype}
\usepackage{algorithm}
\usepackage{algpseudocode}
\usepackage[colorlinks=true,
            linkcolor=blue!50!black,
            urlcolor=blue!50!black,
            citecolor=blue!50!black]{hyperref}

% ---------- Theorem environments ----------
\theoremstyle{plain}
\newtheorem{theorem}{Theorem}[section]
\newtheorem{proposition}[theorem]{Proposition}
\newtheorem{lemma}[theorem]{Lemma}
\newtheorem{corollary}[theorem]{Corollary}

\theoremstyle{definition}
\newtheorem{definition}[theorem]{Definition}

\theoremstyle{remark}
\newtheorem{remark}[theorem]{Remark}
\newtheorem{example}[theorem]{Example}

% ---------- Math shortcuts ----------
\newcommand{\R}{\mathbb{R}}
\newcommand{\N}{\mathbb{N}}
\newcommand{\E}{\mathbb{E}}
\newcommand{\Var}{\operatorname{Var}}
\newcommand{\Cov}{\operatorname{Cov}}
\newcommand{\Corr}{\operatorname{Corr}}
\newcommand{\F}{\mathcal{F}}
\newcommand{\Prob}{\mathbb{P}}
\newcommand{\Q}{\mathbb{Q}}
\newcommand{\dd}{\mathop{}\!\mathrm{d}}
\newcommand{\eqdef}{\overset{\mathrm{def}}{=}}
\newcommand{\indic}{\mathbf{1}}

% ---------- Document metadata ----------
\title{Asian Options by Monte Carlo:\\
       Arithmetic Pricing with the Geometric Control Variate}
\author{Quant Finance Portfolio --- Phase 2, Block 5}
\date{\today}

\begin{document}
\maketitle

\begin{abstract}
Asian options pay on a time-average of the underlying rather than
its terminal value. The arithmetic-average Asian call has no closed
form (the arithmetic average of lognormals is not lognormal), so
Monte Carlo is the standard pricing tool. The geometric-average
Asian call \emph{does} have a closed form, because the geometric
mean of lognormals is lognormal. The two averages are very highly
correlated under typical parameters ($\rho > 0.999$ for a one-year
ATM call with daily monitoring), so the geometric Asian is a
near-perfect control variate for the arithmetic Asian. This block
derives the closed form for the geometric Asian, implements the
three pricers (arithmetic IID, geometric IID, arithmetic with
geometric CV), and demonstrates a variance reduction factor on the
order of $10^3$ at the canonical ATM point.
\end{abstract}

\tableofcontents

\section{Path-dependent payoffs and the Asian family}
\label{sec:setting}

A \emph{path-dependent} option pays at maturity according to the
trajectory of the underlying, not just its terminal value. Define
a discrete monitoring grid $0 = t_0 < t_1 < \cdots < t_N = T$
(here equispaced: $t_k = kT/N$), and let $S_k \eqdef S_{t_k}$.

The two Asian averages are:
\begin{equation}
\bar S \;\eqdef\; \frac{1}{N}\sum_{k=1}^N S_k
\qquad\text{(arithmetic average)},
\end{equation}
\begin{equation}
\tilde S \;\eqdef\; \biggl(\prod_{k=1}^N S_k\biggr)^{1/N}
\qquad\text{(geometric average)}.
\end{equation}
Note that $S_0$ is excluded from both averages (it is the known spot,
not part of the random history).

The arithmetic-average call pays $\max(\bar S - K, 0)$ at $T$; its
risk-neutral price is
\begin{equation}
C^A \;\eqdef\; e^{-rT}\,\E[\max(\bar S - K, 0)].
\end{equation}
The geometric-average call pays $\max(\tilde S - K, 0)$ with price
$C^G$ analogously.

\paragraph{Why arithmetic Asian has no closed form.} Under GBM, each
$S_k$ is lognormal. The product of lognormals is lognormal (the log
is a sum of normals), so $\tilde S$ is lognormal. But the
\emph{sum} of lognormals is not lognormal, so $\bar S$ is not
lognormal. There is no known elementary distribution for $\bar S$,
hence no Black-Scholes formula. Monte Carlo is the workhorse tool.

\paragraph{Why geometric Asian has a closed form.} As noted,
$\log \tilde S = \frac{1}{N}\sum_{k=1}^N \log S_k$, which is normal
because each $\log S_k$ is normal. The first two moments determine
the distribution; we compute them explicitly in
Section~\ref{sec:closed-form}.

\section{Closed form for the geometric Asian call}
\label{sec:closed-form}

Under risk-neutral GBM,
\begin{equation}
S_k \;=\; S_0 \exp\bigl((r - \tfrac{1}{2}\sigma^2)t_k + \sigma W_{t_k}\bigr),
\end{equation}
so
\begin{equation}
\log S_k \;=\; \log S_0 + (r - \tfrac{1}{2}\sigma^2)t_k + \sigma W_{t_k}.
\end{equation}

\subsection{Mean and variance of $\log \tilde S$}

Take the average over $k$:
\begin{equation}
\log \tilde S
\;=\; \log S_0 + (r - \tfrac{1}{2}\sigma^2) \cdot \frac{1}{N}\sum_{k=1}^N t_k
+ \sigma \cdot \frac{1}{N}\sum_{k=1}^N W_{t_k}.
\end{equation}
With equispaced monitoring $t_k = kT/N$:
\begin{equation}
\frac{1}{N}\sum_{k=1}^N t_k
\;=\; \frac{T}{N^2}\sum_{k=1}^N k
\;=\; \frac{T(N+1)}{2N}.
\end{equation}
So the mean is
\begin{equation}
\mu_G \;\eqdef\; \E[\log \tilde S]
\;=\; \log S_0 + (r - \tfrac{1}{2}\sigma^2) \cdot \frac{T(N+1)}{2N}.
\end{equation}

For the variance, we need
\begin{equation}
\Var\!\left(\frac{1}{N}\sum_{k=1}^N W_{t_k}\right)
\;=\; \frac{1}{N^2} \sum_{j=1}^N \sum_{k=1}^N \Cov(W_{t_j}, W_{t_k})
\;=\; \frac{1}{N^2} \sum_{j,k} \min(t_j, t_k).
\end{equation}
Using $t_k = kT/N$:
\begin{equation}
\frac{1}{N^2}\sum_{j,k} \min(t_j, t_k)
\;=\; \frac{T}{N^3}\sum_{j,k}\min(j, k).
\end{equation}
The double sum $\sum_{j,k}\min(j,k)$ over $\{1, \ldots, N\}^2$ equals
\begin{equation}
\sum_{j=1}^N\sum_{k=1}^N \min(j,k)
\;=\; \sum_{m=1}^N m \cdot (2(N-m) + 1)
\;=\; \frac{N(N+1)(2N+1)}{6}.
\end{equation}
(The identity $\sum_{m=1}^N m(2N - 2m + 1) = \frac{N(N+1)(2N+1)}{6}$
is a standard finite sum manipulation.)
So
\begin{equation}
\Var\!\left(\frac{1}{N}\sum W_{t_k}\right)
\;=\; \frac{T(N+1)(2N+1)}{6N^2},
\end{equation}
and
\begin{equation}
\sigma_G^2 \;\eqdef\; \Var(\log \tilde S)
\;=\; \sigma^2 \cdot \frac{T(N+1)(2N+1)}{6N^2}.
\end{equation}

\subsection{Effective parameters and the Black-Scholes formula}

We have $\log \tilde S \sim \mathcal{N}(\mu_G, \sigma_G^2)$, so
$\tilde S$ is lognormal. Define the \emph{effective volatility}
\begin{equation}
\sigma_{\mathrm{eff}}
\;\eqdef\; \sigma \sqrt{\frac{(N+1)(2N+1)}{6N^2}}.
\end{equation}
Note $\sigma_{\mathrm{eff}}^2 \cdot T = \sigma_G^2$. To match the
standard BS form, we want to write
\begin{equation}
\log \tilde S \;=\; \log S_0 + (r_{\mathrm{eff}} - \tfrac{1}{2}\sigma_{\mathrm{eff}}^2)T + \sigma_{\mathrm{eff}}\sqrt T \tilde Z
\end{equation}
for some standard normal $\tilde Z$. Matching means:
\begin{equation}
r_{\mathrm{eff}} \cdot T - \tfrac{1}{2}\sigma_{\mathrm{eff}}^2 T
\;=\; (r - \tfrac{1}{2}\sigma^2) \cdot \frac{T(N+1)}{2N},
\end{equation}
so
\begin{equation}
r_{\mathrm{eff}} \;=\; \frac{(N+1)}{2N}(r - \tfrac{1}{2}\sigma^2)
+ \tfrac{1}{2}\sigma_{\mathrm{eff}}^2.
\end{equation}

Then the geometric Asian call has Black-Scholes form, but with
discounting at $r$ (the original risk-free rate, since we
\emph{discount}, not \emph{drift}, the geometric average):
\begin{equation}
\boxed{
C^G \;=\; e^{-rT}\bigl[ S_0 e^{r_{\mathrm{eff}}\,T}\,\Phi(d_1)
- K\,\Phi(d_2)\bigr]
}
\end{equation}
where
\begin{equation}
d_1 = \frac{\log(S_0/K) + (r_{\mathrm{eff}} + \tfrac{1}{2}\sigma_{\mathrm{eff}}^2)\,T}
{\sigma_{\mathrm{eff}}\sqrt T},
\qquad
d_2 = d_1 - \sigma_{\mathrm{eff}}\sqrt T.
\end{equation}

\paragraph{Sanity check at $N \to \infty$.} As $N$ grows,
\[
\frac{(N+1)(2N+1)}{6N^2} \;\to\; \frac{1}{3},
\]
so $\sigma_{\mathrm{eff}} \to \sigma/\sqrt 3$. This reproduces the
classical result that the continuous-monitoring geometric Asian
behaves as a vanilla call with volatility $\sigma/\sqrt 3$. Good
sanity check.

\paragraph{Sanity check at $N = 1$.} With $N = 1$ the geometric
average reduces to $S_T$ itself, so $C^G$ should reduce to the
vanilla call price. Plugging in:
\[
\sigma_{\mathrm{eff}}^2 = \sigma^2 \cdot \frac{2 \cdot 3}{6} = \sigma^2,
\quad
r_{\mathrm{eff}} = (r - \tfrac{1}{2}\sigma^2) \cdot 1 + \tfrac{1}{2}\sigma^2 = r,
\]
so $C^G = e^{-rT}[S_0 e^{rT}\Phi(d_1) - K\Phi(d_2)] = S_0\Phi(d_1) -
K e^{-rT}\Phi(d_2)$, the standard BS formula. Verified.

\section{Why geometric is a near-perfect control for arithmetic}
\label{sec:correlation}

The geometric and arithmetic means satisfy the AM-GM inequality:
$\tilde S \le \bar S$ pointwise. Moreover, on a single GBM path,
both quantities respond to the same randomness (the same Brownian
increments), so they are highly correlated. We can quantify how
high.

\paragraph{Heuristic argument.} Write $X_k = \log S_k$. Then
\[
\log \tilde S = \frac{1}{N}\sum_k X_k,
\qquad
\log \bar S = \log\!\left(\frac{1}{N}\sum_k e^{X_k}\right).
\]
For small relative variation in $\{X_k\}$ (which holds for
moderate $\sigma\sqrt T$ and reasonable $N$), $\log \bar S \approx
\log \tilde S + \tfrac{1}{2}\Var_k(X_k)$ by a Taylor expansion of
the log-sum-exp around its mean. The correction is a deterministic
shift, so $\Corr(\log \tilde S, \log \bar S) \to 1$ in this regime.

In practice, for a 1-year ATM call with $\sigma = 0.20$ and $N = 50$
monitoring dates, empirical correlation between $\bar S - K$ and
$\tilde S - K$ exceeds $0.9999$ on $10^4$ Monte Carlo paths.

\paragraph{Implication for variance reduction.} The control variate
identity from Block~2.2 gives, for paired estimators
$X = e^{-rT}\max(\bar S - K, 0)$ and $Y = e^{-rT}\max(\tilde S - K, 0)$,
\begin{equation}
\Var(X - \beta^\star Y) \;=\; \Var(X)(1 - \rho^2),
\end{equation}
where $\beta^\star = \Cov(X, Y)/\Var(Y)$ and $\rho = \Corr(X, Y)$.
With $\rho \approx 0.9999$, the variance reduction factor is
$1/(1 - \rho^2) \approx 5\,000$. This is by far the most powerful
control variate we have encountered in Phase~2.

\section{Implementation: three pricers}
\label{sec:impl}

The block ships three pricers in a new module
\texttt{quantlib.asian} (and the C++ mirror).

\subsection{The arithmetic IID baseline}

Standard MC: simulate $n$ paths of $\{S_k\}_{k=1}^N$ exactly under
GBM (no Euler), compute $\bar S^{(i)}$ for each, return
\begin{equation}
\hat C^A_n \;=\; \frac{1}{n}\sum_{i=1}^n e^{-rT}\max(\bar S^{(i)} - K, 0).
\end{equation}

\begin{algorithm}[H]
\caption{Arithmetic Asian call by IID Monte Carlo (exact GBM)}
\begin{algorithmic}[1]
\Require Spot $S_0$, params $(K, r, \sigma, T)$, monitoring $N$, sample $n$.
\State Compute $h \gets T/N$, drift $\mu \gets (r - \tfrac{1}{2}\sigma^2)h$, diffusion $\nu \gets \sigma\sqrt h$.
\For{$i = 1, \ldots, n$}
\State $S \gets S_0$
\State $\text{sum} \gets 0$
\For{$k = 1, \ldots, N$}
\State $Z \gets \mathcal{N}(0, 1)$
\State $S \gets S \cdot \exp(\mu + \nu Z)$
\State $\text{sum} \gets \text{sum} + S$
\EndFor
\State $\bar S^{(i)} \gets \text{sum}/N$
\State $\Pi^{A,(i)} \gets e^{-rT}\max(\bar S^{(i)} - K, 0)$
\EndFor
\State \Return mean and half-width of $\{\Pi^{A,(i)}\}$.
\end{algorithmic}
\end{algorithm}

\subsection{The geometric IID estimator (validation)}

Same path simulation, but accumulate $\sum \log S_k$ and exponentiate
at the end. Used purely to validate the closed form: the IID estimate
must agree with $C^G$ from Section~\ref{sec:closed-form} within MC
error.

\subsection{The arithmetic with geometric CV}

Run both estimators on the \emph{same paths} (CRN), compute the
optimal $\beta^\star$ from sample covariance, and apply the control
variate correction. Following Block~2.2's API:
\begin{equation}
\hat C^A_{\mathrm{CV}}
\;=\; \bar X - \hat\beta^\star (\bar Y - C^G),
\end{equation}
where $C^G$ is the closed-form geometric price. The half-width is
computed from the variance of $X - \hat\beta^\star Y$, which is
$(1 - \hat\rho^2)$ times the variance of $X$.

\section{Empirical validation strategy}
\label{sec:validation}

The script \texttt{python/validate\_mc\_asian.py} runs four tests.

\paragraph{Test 1: closed-form match for geometric Asian.} The IID
estimator $\hat C^G_n$ at $n = 10^5$ should agree with the closed
form $C^G$ within $3$ half-widths.

\paragraph{Test 2: CV gives huge VRF.} Compare half-width of
$\hat C^A_{\mathrm{CV}}$ vs $\hat C^A_{\mathrm{IID}}$ at the same
budget. Predicted: VRF $> 500$ (typical $> 1000$ at our parameters).

\paragraph{Test 3: empirical correlation.} Measure $\hat \rho$
between arithmetic and geometric per-path payoffs. Predicted:
$\rho > 0.999$.

\paragraph{Test 4: input validation.} Mechanical.

\section{Contrast with Block~2.2 control variates}
\label{sec:contrast}

Block~2.2 introduced control variates with two examples on the
European call: discounted-underlying ($\rho \approx 0.92$, VRF $\approx 7$)
and asset-or-nothing ($\rho \approx 0.77$, VRF $\approx 2.5$). Block~5
applies the same machinery to a different problem and gets a vastly
better outcome ($\rho > 0.999$, VRF $> 1000$).

The lesson is that \emph{control variate quality is determined by
the structural relationship between control and target}, not by the
sophistication of the technique. The geometric and arithmetic
averages of the same set of lognormals are bound by AM-GM and share
all the same noise; they could not help being correlated. The
discounted-underlying and the European call payoff are connected
only through the indicator function $\indic_{\{S_T > K\}}$, which
loses about $40\%$ of the variance to the bounded-payoff truncation.

\paragraph{Operational rule.} When designing a control variate, look
for a problem that shares as much algebraic structure as possible
with the target. The Asian-geometric-as-CV-for-Asian-arithmetic is
the canonical example of this principle.

\section{Bibliographic notes}
\label{sec:biblio}

The geometric-Asian closed form has been re-derived independently
many times; the version we use is essentially that of Kemna and
Vorst (1990)~\cite{KemnaVorst1990}. The use of the geometric Asian
as a control variate for the arithmetic Asian goes back to the same
paper. Glasserman~\cite[Section 4.5]{Glasserman} treats the
construction at modern length and discusses the limits of the
approximation. Asian options are also treated extensively in
Hull~\cite{Hull}.

\begin{thebibliography}{99}

\bibitem{KemnaVorst1990}
A.\ G.\ Z.\ Kemna and A.\ C.\ F.\ Vorst,
\emph{A pricing method for options based on average asset values},
Journal of Banking and Finance, 14(1):113--129, 1990.

\bibitem{Glasserman}
P.\ Glasserman,
\emph{Monte Carlo Methods in Financial Engineering},
Springer, 2004.

\bibitem{Hull}
J.\ C.\ Hull,
\emph{Options, Futures, and Other Derivatives},
10th edition, Pearson, 2018.

\end{thebibliography}

\end{document}
